Todos los días nos vemos en un espejo. Casi nunca nos preguntamos dónde está, exactamente, la persona que vemos del otro lado.

Parece una pregunta tonta, pero no lo es.

Cuando nos paramos frente a un espejo, la luz que sale de nuestra cara rebota en el vidrio y llega a nuestros ojos, y el cerebro, que siempre supone que la luz viaja en línea recta, prolonga esos rayos hacia atrás hasta un punto donde parecen haber nacido, un punto que queda detrás del espejo, donde en realidad no hay nada. Por eso a esa imagen se le llama virtual: se puede ver, pero no se puede proyectar en una pantalla.

En las dos prácticas anteriores trabajamos con la ley de la reflexión y con la refracción. Esta vez la idea fue aplicar la primera a un caso muy cotidiano, el del espejo plano, colocando objetos a distintas distancias de un espejo y midiendo dónde aparecía su imagen para comprobar si de verdad queda a la misma distancia detrás que el objeto delante.
